% GENERATED FILE. Edit canonical lessons, then run npm run generate:books.
% canonical-source-hash: fnv1a64:83c9122d2a6a8b58
% canonical-lessons: TE-C150-vangmayam, TE-S171-letter-nga

\chapter{A Word for Literature — Then the Letter Inside It}
\label{ch:te-vangmayam-nga}
\hlchaptermodality{150}

\begin{chapteropening}
\textbf{By the end of this chapter:} \emph{I can say the Telugu word for literature, recognise the rare letter inside it, and write that letter in a source-backed order.}

Take \te{ఙ} from the already-known word \te{వాఙ్మయం} and write its five movements in three pen-down runs.
\end{chapteropening}

\section[\texorpdfstring{\te{వాఙ్మయం} (vāṅmayaṁ)}{వాఙ్మయం (vāṅmayaṁ)}]{\te{వాఙ్మయం} (vāṅmayaṁ) — literature, the whole body of expression in words}

\label{lesson:TE-C150-vangmayam}



\begin{quote}
You already know \textbf{\te{పుస్తకం}}, a book, and you have held a short passage without translating every word. This lesson names the larger body that books and spoken works belong to.
\end{quote}

\subsection*{You'll want to know: \te{వాఙ్మయం}}
\textbf{\te{వాఙ్మయం}} — \emph{vāṅmayaṁ} — \textquotedblleft{}literature\textquotedblright{}, or more broadly, expression made of words.

Say it in three gentle pieces: \emph{vāṅ · ma · yaṁ}. The first piece ends with the velar nasal heard in English \emph{sing}. The spelling keeps that sound as \textbf{\te{ఙ్}} before \textbf{\te{మ}}: \textbf{\te{వా}-\te{ఙ్}-\te{మ}-\te{యం}}.

The word came through Sanskrit: \emph{vāc} is speech, and \emph{maya} means made of or filled with. A single \textbf{\te{పుస్తకం}} is one book; \textbf{\te{వాఙ్మయం}} is the world of words that books, speeches, and other verbal works make together.

\begin{grammarlens}[title={What you've built}]
One book, then the wider field it can belong to. The unusual spelling has had a meaning before it becomes a handwriting exercise.
\end{grammarlens}

\subsection*{Your turn}
\begin{itemize}
\raggedright
  \item \emph{Say it:} \emph{vāṅ · ma · yaṁ}, slowly
  \item \emph{Say it:} \emph{vāṅmayaṁ}, as one word
  \item \emph{Read it:} \textbf{\te{పుస్తకం} · \te{వాఙ్మయం}} — one book, then literature
  \item \emph{Find:} the small \textbf{\te{ఙ్}} immediately before \textbf{\te{మ}}
\end{itemize}

\begin{tcolorbox}[breakable,colback=teal!4,colframe=teal!35!black,title={Before you move on}]
What does \textbf{\te{వాఙ్మయం}} mean? (\textquotedblleft{}Literature; the wider body of expression in words.\textquotedblright{}) Say it once more.
\end{tcolorbox}
\section[\texorpdfstring{\te{ఙ} (ṅa)}{ఙ (ṅa)}]{\te{ఙ} — the quiet letter inside \te{వాఙ్మయం}}

\label{lesson:TE-S171-letter-nga}



\begin{quote}
Say \textbf{\te{వాఙ్మయం}} once. What does it mean? (\textquotedblleft{}Literature; the wider body of expression in words.\textquotedblright{})

Now slow the first two pieces: \emph{vāṅ · ma}. The letter between them is today's whole task.
\end{quote}

\subsection*{Script you'll notice: \te{ఙ}}
\textbf{\te{ఙ}} — \emph{ṅa}.

It is the velar nasal: the sound at the end of English \emph{sing}. In \textbf{\te{వాఙ్మయం}}, the vowel is suppressed, so \textbf{\te{ఙ్}} leans into the following \textbf{\te{మ}}. The base letter you will write is \textbf{\te{ఙ}}.

This is an uncommon letter in modern Telugu words. That is why the book waited until it had given you a real word that needs it.

\subsection*{Writing: \te{ఙ} — five small movements}
\hlblockfigure{\detokenize{figures/TE-S171-letter-nga-filmstrip.pdf}}{How \te{ఙ} is written, stroke by stroke}

Follow the filmstrip slowly. Movements 1–3 stay joined: turn around the compact upper-left lobe, continue around the broad lower bowl, then curl up through the right lobe. Lift for the inner horizontal bar. Lift once more for the short downward headstroke.

The numbered route is one attested school-style order; Telugu handwriting can vary. Keep the three pen-down runs clear rather than racing the shape.

\subsection*{Your turn}
\begin{itemize}
\raggedright
  \item \emph{Trace it:} \textbf{\te{ఙ}} once, following the five numbered movements
  \item \emph{Copy:} \textbf{\te{ఙ}} twice without tracing
  \item \emph{Find:} \textbf{\te{ఙ్}} inside \textbf{\te{వాఙ్మయం}}
  \item \emph{Say it:} \emph{vāṅmayaṁ} as you finish
\end{itemize}

\begin{tcolorbox}[breakable,colback=teal!4,colframe=teal!35!black,title={Before you move on}]
Which letter is this — \textbf{\te{ఙ}}? What sound does it carry? (\emph{ṅa}.) Which word gave it to you? (\textbf{\te{వాఙ్మయం}}.)
\end{tcolorbox}
